\section{Overview}


\textit{Normative language per \href{https://tools.ietf.org/html/rfc2119}{RFC~2119}.}

\subsection{Scope}

This specification governs the integration of a tenant's Single Sign-On (SSO)
Identity Provider (IdP) with Canvasflow. It covers the Progressive Web Application (PWA)
and its backing GraphQL API server, defines the requirements a tenant MUST satisfy to
issue access tokens that the Canvasflow Resource Server will accept, the entitlement
model by which content access is granted, the onboarding process and mandatory
validation gate, and the guarantees Canvasflow provides in return.

This specification is version 1.0.0 and is the single canonical reference for 
all new integrations.

\subsection{Division of Trust}
The architecture establishes a clean, auditable boundary between two parties, each
responsible for exactly the functions they own. This boundary follows the role
definitions of \href{https://datatracker.ietf.org/doc/html/rfc6749#section-1.1}{RFC~6749~\S1.1.}

\textbf{The tenant operates the Authorization Server.} All authentication decisions
(credential verification, MFA, session management, social connections) happen entirely
within the tenant's IdP. All authorization decisions --- which users hold which
entitlements, when a session is created or revoked --- are the tenant's responsibility.
The tenant signs every token with keys or secrets they control. Canvasflow never
participates in authentication and never inspects a user's credentials.


\textbf{Canvasflow operates a pure Resource Server} \href{https://datatracker.ietf.org/doc/html/rfc6749#section-1.1}{\textbf{RFC~6749~\S1.1.}} The Canvasflow
GraphQL server does exactly two things: (a)~validate the cryptographic signature and
required claims of every incoming JWT access token, and (b)~return content gated by the
verified entitlement claim. It does not revoke or invalidate tokens, does not maintain
a token blocklist, does not call the tenant IdP at request time (JWKS fetching is the
only exception, and it uses a server-side cache), makes no authorization decision of its
own --- it only \emph{enforces} decisions the tenant's Authorization Server already sealed
into the signed token, and stores no user accounts, passwords, or sessions.

The signed JWT access token is the sole interface between the two systems. Its signature
is the trust anchor. The Resource Server needs nothing else.

\subsection{The Flow in One Sentence}

A tenant's user authenticates via OAuth~2.0 Authorization Code flow with PKCE \linebreak
(\href{https://datatracker.ietf.org/doc/html/rfc6749#section-4.1}{RFC~6749~\S4.1} + 
\href{https://datatracker.ietf.org/doc/html/rfc7636}{RFC~7636}), receiving a JWT access token 
\href{https://datatracker.ietf.org/doc/html/rfc9068}{RFC~9068} whose entitlement
claim was pre-populated by the tenant's IdP hook with Canvasflow resource identifiers
mapped from the user's internal entitlements before signing --- and the Canvasflow
GraphQL server validates that token's signature and claims on every request, gating
content strictly by the verified entitlement set without any further call to the tenant.

\pagebreak

\subsection{Out of Scope}\label{sec:out-of-scope}

The following items are explicitly out of scope for this specification and for the
Canvasflow Resource Server in its current version:

\begin{itemize}
  \item \textbf{Token revocation.} Canvasflow has no revocation endpoint, no token
    blocklist, and no introspection endpoint. Revocation is the Authorization Server's
    responsibility. The maximum exposure window for a revoked access is bounded by the
    token's \texttt{exp} claim, which MUST NOT exceed 3600~seconds.

  \item \textbf{Server-side entitlement mapping.} In this version, the
    tenant owns and executes all mapping between their internal identifiers and
    Canvasflow resource IDs. Canvasflow never performs this mapping at runtime.

  \item \textbf{User provisioning and SCIM sync.} Canvasflow does not accept SCIM
    provisioning requests and does not maintain a user directory of tenant users.

  \item \textbf{Authorization logic on the Canvasflow side.} The Canvasflow server
    enforces entitlements carried in the token. It invents no additional access rules.

  \item \textbf{Multi-factor authentication (MFA) enforcement.} MFA is the tenant's
    domain. Canvasflow has no visibility into or control over the tenant's authentication
    strength requirements.
\end{itemize}

\pagebreak

\subsection{Architecture Sequence Diagrams}

The following diagrams document the three phases of the architecture: the 
one-time onboarding setup, the per-user authentication flow, and the 
per-request validation flow. All claim names and value formats reflect 
the requirements of this specification(\hyperref[sec:entitlement-mapping]{\S4}).

\subsubsection{Onboarding (One-Time, Out-of-Band)}

\begin{figure}[ht!]
    \centering % Centers the image horizontally
    \includegraphics[width=1\textwidth]{on-boarding-diagram.png} % Injects and resizes the image
    \caption{Tenant Onboarding Sequence Diagram} % Caption for the image
    \label{fig:on-boarding} % Used for cross-referencing
\end{figure}

\pagebreak

\subsubsection{Authentication Time (Per-User Login)}

\begin{figure}[ht!]
    \centering % Centers the image horizontally
    \includegraphics[width=1\textwidth]{authentication-diagram.png} % Injects and resizes the image
    \caption{Tenant Authentication Sequence Diagram} % Caption for the image
    \label{fig:authentication} % Used for cross-referencing
\end{figure}

\pagebreak

\subsubsection{Request Time (Per-API Call)}

\begin{figure}[ht!]
    \centering % Centers the image horizontally
    \includegraphics[width=1\textwidth]{request-diagram.png} % Injects and resizes the image
    \caption{Tenant Request Sequence Diagram} % Caption for the image
    \label{fig:request} % Used for cross-referencing
\end{figure}